%%%%%%%%%%  scan idx 47  |  folio 34  %%%%%%%%%%
\emph{fait opérer le groupe symétrique $\mathfrak S_p$ sur $N^p$, donc sur $N^{(p)}$ donc sur $\Lambda(N^{(p)})$, puis sur $\bigotimes^p F$ et sur $\Lambda(N^{(p)})\otimes\bigotimes^p F$. Considérons cet espace vectoriel comme un module gradué sur $\mathfrak S_p\times\Gl(N)\times\Gl(F)$ et prenons sa «~partie alternée~» qui est un module gradué sur $\Gl(N)\times\Gl(F)\eqqcolon G$. On a alors}
\begin{equation}\tag{1.33}
\lambda^p(N,F)=\underline{E}_G\left(\left(\Lambda(N^{(p)})\otimes\bigotimes^p F\right)^{\mathrm{alt}}\right).
\end{equation}

N.B. Si $M$ est un complexe de $G$-modules, on note $\underline{E}_G(M)$ la somme alternée dans $\underline{K}(G)$ des classes des composantes de $M$~; c'est la caractéristique d'Euler--Poincaré de $M$~; dans la formule (1.33), le premier membre $\lambda^p(N,F)$ désigne l'élément de l'anneau $\underline{K}(G)$ obtenu en substituant aux variables $\lambda^iN$ et $\lambda^iF$ les classes dans $\underline{K}(G)$ des $G$-modules $\Lambda^i(N)$ et $\Lambda^i(F)$ respectivement.

\textbf{Démonstration :} En développant le deuxième membre de (1.32) par la deuxième formule (1.11), on voit que la première assertion est prouvée si on prouve l'assertion analogue obtenue en faisant $x=1$, i.e. si l'on prouve que dans le $\lambda$-anneau $\Z[\lambda^1N,\ldots,\lambda^qN]$ où $\lambda^iN=0$ pour $i>q$, l'élément $\lambda^p(\lambda_{-1}(N))$ est divisible par $\lambda_{-1}(N)$. A priori, le quotient $\lambda^p(N,1)=\lambda^p(\lambda_{-1}(N))/\lambda_{-1}(N)$ est une série formelle contenue dans le corps des quotients $K$ de l'anneau des séries formelles en les $\lambda^iN$ pour $1\leq i\leq q$. D'autre part, désignons aussi par $N$ un vectoriel de dimension $q$ sur un corps algébriquement clos $k$ de caractéristique $0$, et considérons l'élément $\ell^p(N)=\underline{E}_G((\Lambda(N^{(p)})^{\mathrm{alt}})\in \underline{K}(G)$, où maintenant l'on a $G=\Gl(N)$. On sait (\S 2, théorème 1.3) que $\underline{K}(G)$ s'identifie 
